% Generated draft: requires mathematical and source review.
\documentclass[11pt]{article}
\usepackage[margin=25mm]{geometry}
\usepackage{fontspec}
\setmainfont{Latin Modern Roman}
\usepackage{amsmath,amssymb,mathtools,unicode-math}
\setmathfont{Latin Modern Math}
\usepackage{enumitem,needspace,xurl,hyperref}
\setlength{\parindent}{0pt}
\setlength{\parskip}{5pt}
\setlength{\emergencystretch}{3em}
\newcommand{\N}{\mathbb{N}}
\newcommand{\Z}{\mathbb{Z}}
\newcommand{\Q}{\mathbb{Q}}
\newcommand{\R}{\mathbb{R}}
\newcommand{\C}{\mathbb{C}}
\newcommand{\F}{\mathbb{F}}
\newcommand{\A}{\mathbb{A}}
\newcommand{\PP}{\mathbb{P}}
\newcommand{\E}{\mathbb{E}}
\newcommand{\eps}{\varepsilon}
\newcommand{\dd}{\,\mathrm{d}}
\begin{document}
\section*{UnsolvedMath ID: 30000754}
\textbf{Source problem code:} catalogue entry
\subsection{Definitions and mathematical statement}
A real number $x$ is a Liouville number if it is irrational and \[\forall m\in\Z_{\geq1}\ \exists p\in\Z\ \exists q\in\Z_{\geq2}:\qquad 0<\left|x-\frac pq\right|<q^{-m}.\] Let $e$ be the base of the natural logarithm and let $\pi$ be the usual circle constant. The source asks whether $e^\pi$ is a Liouville number and explicitly conjectures the negative answer: \[e^\pi\text{ is not a Liouville number}.\] In view of the transcendence reported in the supplied literature assessment, the conjecture may equivalently be written \[\exists m\in\Z_{\geq1}\ \forall p\in\Z\ \forall q\in\Z_{\geq2}:\qquad \left|e^\pi-\frac pq\right|\geq q^{-m}.\] The source's literature triage, dated 2026-08-21, labels the classification open and reports known transcendence and quantitative approximation bounds. It does not state those bounds, and they are not treated here as establishing the conjectured non-Liouville property.
\subsection{Short English statement}
Determine whether $e^\pi$ admits rational approximations of arbitrarily large power-law accuracy. The report conjectures that it does not, and the supplied 2026-08-21 assessment labels this open.
\subsection{Sources}
\begin{itemize}\item Local UnsolvedMath registry, permanent ID 30000754, source problem catalogue entry, \emph{Liouville Classification of the Gelfond Constant}, proposed year 2007; no individual proposer supplied.\item Supplied report citation: \emph{Diophantische Approximationen} (2008), Oberwolfach Reports, DOI 10.4171/owr/2007/21, pp.\ 1115--1190, \url{https://doi.org/10.4171/owr/2007/21}.\item Supplied literature label: Oberwolfach report passage, \url{https://ems.press/content/serial-article-files/46108}.\item Supplied literature label: Waldschmidt, Schanuel conjecture notes, \url{https://webusers.imj-prg.fr/~michel.waldschmidt/articles/pdf/SchanuelEn.pdf}.\end{itemize}
\end{document}
